\section{Результаты экспериментов и их интерпретация}
\label{sec:experiment-results}

\subsection{Полная выгрузка}

Измерения показывают, что вариант хранилища, основанный на
\texttt{ranger}, быстрее выполняет полную выгрузку таблицы. Это объясняется
тем, что он концептуально хранит уже подготовленный линейный набор
записей и может отдавать его почти как копию готового массива.
Шардированный ART, напротив, должен обойти все шарды, декодировать
ключи, собрать raw-префиксы и построить непересекающуюся таблицу для
внешней выдачи.

Конкретные численные результаты для $n=100000$ приведены в
табл.~\ref{tab:bench-dump}. Во всех ART-сценариях сортировка дампа
отключена; измерение включает построение непересекающейся таблицы.

\begin{table}[ht]
\centering
\caption{Бенчмарк \texttt{Dump} для $n=100000$}
\label{tab:bench-dump}
\footnotesize
\begin{tabular}{>{\raggedright\arraybackslash}p{4.2cm}rr}
\toprule
Сценарий & \texttt{ranger} & ART \\
\midrule
IPv4 \texttt{Dump} & $48{,}90$ мс/оп. & $142{,}80$ мс/оп. \\
IPv6 \texttt{Dump} & $101{,}36$ мс/оп. & $201{,}93$ мс/оп. \\
\bottomrule
\end{tabular}
\end{table}

Эти данные фиксируют проигрыш ART по полной выгрузке. Для целевого профиля
нагрузки это не опровергает гипотезу, поскольку \texttt{Dump} не является
основным рабочим режимом потоковой модели, но задаёт ограничение для
сценариев частых полных выгрузок.

\subsection{Инкрементальное обновление}

С точки зрения постановки задачи этот раздел -- центральный: он
показывает, что потоковое применение одиночных BGP-изменений к
in-memory базе технически осуществимо, а не только объявлено в
архитектуре. Именно на операции одиночного \texttt{Upsert} ART-подход
демонстрирует решающее преимущество (табл.~\ref{tab:bench-upsert}).
Обе реализации поддерживают точечное изменение состояния, но ART
выполняет эту операцию локально в одном шарде и требует меньше времени и
аллокаций.

\begin{table}[ht]
\centering
\caption{Бенчмарк одиночного \texttt{Upsert} поверх заполненной
таблицы, $n=100000$}
\label{tab:bench-upsert}
\small
\begin{tabular}{lrr}
\toprule
Сценарий & \texttt{ranger} & ART \\
\midrule
IPv4 \texttt{Upsert} & $3{,}94$ мкс/оп. & $173{,}1$ нс/оп. \\
IPv6 \texttt{Upsert} & $6{,}43$ мкс/оп. & $199{,}2$ нс/оп. \\
\bottomrule
\end{tabular}
\end{table}

Разница составляет десятки раз при $n=10^5$. Это главный эмпирический
аргумент в пользу ART для потоковой модели: отдельное BGP-событие
применяется с малой задержкой и небольшим числом аллокаций, а значит
инкремент становится естественным базовым режимом внутри выделенного
микросервиса с общим API.

\subsection{Lookup по IP}

Это вспомогательный разрез относительно центральной цели работы
(микросервис и поток обновлений), но он фиксирует цену точечного чтения
после выбора структуры. По операции \texttt{Lookup} обе реализации дают
сопоставимую задержку. ART проверяет только активные длины префиксов,
поэтому стоимость поиска не определяется полным диапазоном возможных
масок IPv4 или IPv6. Результаты приведены в табл.~\ref{tab:bench-lookup}.

\begin{table}[ht]
\centering
\caption{Бенчмарк \texttt{Lookup}, $n \approx 10^5$}
\label{tab:bench-lookup}
\small
\begin{tabular}{lrr}
\toprule
Сценарий       & \texttt{ranger}       & ART (16 шардов) \\
\midrule
IPv4 & $\approx 150$--$180\,$нс/оп. & $\approx 150\,$нс/оп. \\
IPv6 & $175{,}5\,$нс/оп. & $211{,}7\,$нс/оп. \\
\bottomrule
\end{tabular}
\end{table}

Эти данные показывают, что \texttt{Lookup} остаётся пригодным для
синхронного gRPC-запроса. ART не является универсальным победителем по
всем метрикам, но после учёта активных длин префиксов разрыв по чтению
практически устранён.

\subsection{Влияние числа шардов}

Эксперименты показывают, что увеличение числа шардов не даёт линейного
ускорения всех операций. Для последовательного \texttt{Lookup} выигрыш
находится в пределах шума. Для \texttt{Dump} большое число шардов
ухудшает результат из-за роста служебных издержек и последующего
глобального flatten (табл.~\ref{tab:bench-shards}). Для конкурентных
update-сценариев шардирование остаётся полезным, потому что снижает
конкуренцию между независимыми модификациями.

\begin{table}[ht]
\centering
\caption{Измерение ART при разном числе шардов на смешанном наборе
$3\cdot10^6$ записей}
\label{tab:bench-shards}
\small
\begin{tabular}{lrr}
\toprule
$K$ (шардов) & \texttt{Lookup} & \texttt{Dump} \\
\midrule
$1$  & $175{,}1$ нс/оп. & $3{,}62$ с/оп. \\
$8$  & $176{,}4$ нс/оп. & $4{,}99$ с/оп. \\
$16$ & $175{,}9$ нс/оп. & $5{,}59$ с/оп. \\
$32$ & $175{,}6$ нс/оп. & $6{,}24$ с/оп. \\
$64$ & $177{,}8$ нс/оп. & $6{,}89$ с/оп. \\
\bottomrule
\end{tabular}
\end{table}

Следовательно, число шардов является параметром настройки, а не
правилом «чем больше, тем лучше». Для dump-heavy сценариев лучший
результат показывает один шард; для конкурентных read/write-нагрузок
практически разумной остаётся зона умеренного шардирования.

\subsection{Поведение IPv6}

IPv6 заслуживает отдельного внимания. По \texttt{Dump} ART остаётся
медленнее \texttt{ranger}, поскольку полная выдача включает построение
непересекающейся таблицы. По \texttt{Upsert} ART сохраняет преимущество:
$199{,}2$ нс/оп. против $6{,}43$ мкс/оп. у \texttt{ranger}. По
\texttt{Lookup} результат сопоставим: $211{,}7$ нс/оп. у ART против
$175{,}5$ нс/оп. у \texttt{ranger}.

\subsection{Смешанный набор на 3 млн записей}

Для проверки масштабирования был отдельно выполнен тяжёлый прогон на
$3\,000\,000$ префиксных записей в смешанном режиме: IPv4 \texttt{/24} и
IPv6 \texttt{/64}. ART использовался с 16 шардами и выключенной
сортировкой полной выгрузки. Результаты приведены в
табл.~\ref{tab:bench-3m-mixed}.

\begin{table}[ht]
\centering
\caption{Сравнение \texttt{ranger} и шардированного ART на смешанном наборе из $3\cdot10^6$ записей}
\label{tab:bench-3m-mixed}
\small
\resizebox{\textwidth}{!}{%
\begin{tabular}{lrrr}
\toprule
Сценарий & \texttt{ranger} & ART (16 шардов, без сортировки) & Победитель \\
\midrule
\texttt{Replace} / полная загрузка
  & $30{,}84$ с/оп.
  & $693{,}98$ мс/оп.
  & ART \\
\texttt{Lookup}
  & $187{,}8$ нс/оп.
  & $176{,}7$ нс/оп.
  & ART \\
\texttt{Dump}
  & $4{,}25$ с/оп.
  & $6{,}16$ с/оп.
  & \texttt{ranger} \\
\texttt{Upsert}
  & $7{,}99$ мкс/оп.
  & $255{,}5$ нс/оп.
  & ART \\
\bottomrule
\end{tabular}
}
\end{table}

Та же серия измерений показывает высокую стоимость полной выгрузки по
памяти. Для \texttt{Dump} на $3\cdot10^6$ записей без искусственного
memory pressure \texttt{ranger} выделяет около $10{,}88$~GiB/оп.,
ART с одним шардом -- около $11{,}41$~GiB/оп., ART с 32 шардами --
около $11{,}82$~GiB/оп. При \texttt{GOMEMLIMIT=1GiB} операция
завершается, но заметно замедляется из-за GC: \texttt{ranger} занимает
около $11{,}09$~с, ART с одним шардом -- около $13{,}71$~с, ART с
8 шардами -- около $11{,}26$~с. Для ART при \texttt{GOMEMLIMIT=2GiB}
штраф в этом бенчмарке почти исчезает.

Этот прогон уточняет общую картину. На большой смешанной таблице ART
выигрывает в полной загрузке состояния, точечном поиске и
инкрементальном обновлении. Основная дорогая операция для ART --
полный \texttt{Dump} с построением непересекающейся таблицы.
Следовательно, выбор хранилища зависит от профиля нагрузки: для
состояния, обновляемого потоком BGP-событий, критичны \texttt{Replace}
и \texttt{Upsert}, тогда как для сценариев частой полной выгрузки нужен
кэшированный или потоковый вариант \texttt{Dump}.

\subsection{Сводный итог сравнения}

Сводка по сценариям приведена в табл.~\ref{tab:bench-summary}.

\begin{table}[ht]
\centering
\caption{Итоговая сводка: какой вариант хранилища выигрывает в каком сценарии}
\label{tab:bench-summary}
\small
\begin{tabular}{>{\raggedright\arraybackslash}p{4.0cm}lp{6.5cm}}
\toprule
Сценарий & Победитель & Комментарий \\
\midrule
Полная выгрузка
  & \texttt{ranger}
  & ART выполняет обход дерева и строит непересекающуюся таблицу \\
Одиночный \texttt{Upsert}
  & ART
  & десятки раз преимущества; решающий фактор для потокового режима \\
\texttt{Lookup} IPv4
  & сопоставимо
  & фильтрация по активным длинам префикса даёт близкую задержку у обеих реализаций \\
\texttt{Lookup} IPv6
  & сопоставимо
  & \texttt{ranger} быстрее примерно в 1{,}2 раза на $n=10^5$ \\
Полная загрузка \texttt{Replace} на $3\cdot10^6$ записей
  & ART
  & $693{,}98$ мс/оп. против $30{,}84$ с/оп. у \texttt{ranger} \\
Потоковое обновление
  & ART
  & отдельное BGP-событие применяется с малой задержкой и небольшим числом аллокаций \\
Потоковое обновление на $3\cdot10^6$ записей
  & ART
  & \texttt{Upsert}: $255{,}5$ нс/оп. против $7{,}99$ мкс/оп. у \texttt{ranger} \\
\bottomrule
\end{tabular}
\end{table}

Из сводки следуют четыре вывода, непосредственно связанные с измеренными
операциями:

\begin{enumerate}
\item \texttt{ranger} сохраняет преимущество в сценариях частой полной
выгрузки: на смешанном наборе $3\cdot10^6$ записей \texttt{Dump} у
\texttt{ranger} занимает $4{,}25$ с/оп. против $6{,}16$ с/оп. у ART;
\item ART лучше соответствует сценарию потокового применения BGP-событий:
одиночный \texttt{Upsert} на смешанном наборе выполняется за
$255{,}5$ нс/оп., тогда как для \texttt{ranger} та же операция занимает
$7{,}99$ мкс/оп.;
\item на полной загрузке большого смешанного набора ART также показывает
меньшее время: \texttt{Replace} занимает $693{,}98$ мс/оп. против
$30{,}84$ с/оп. у \texttt{ranger};
\item практический выбор реализации зависит от профиля нагрузки: для
сервиса, который должен регулярно применять инкрементальные BGP-обновления
и быстро выполнять полную загрузку состояния, измерения подтверждают
преимущество ART; для частых полных выгрузок нужен кэшированный или
потоковый \texttt{Dump}.
\end{enumerate}

Последний вывод задаёт границу применимости результата. Выделение функции
в микросервис и единый gRPC контракт для нескольких продуктов образуют
самостоятельный архитектурный слой; шардированный ART снижает время
потокового обновления по сравнению с \texttt{ranger}-backend-ом.

\subsection{Совокупная стоимость потокового профиля нагрузки}

Чтобы перевести бенчмарки в эксплуатационно содержательную плоскость для
уже выделенного микросервиса, полезно построить простую модель
совокупной стоимости. Для поддержания актуальной таблицы
\texttt{prefix -> ASN} ключевым режимом является не полная выгрузка
состояния, а регулярное применение инкрементальных BGP-обновлений.
Поэтому проигрыш ART в операции \texttt{Dump} сам по себе не определяет
выбор хранилища: его нужно сопоставлять со стоимостью \texttt{Upsert},
\texttt{Delete} и \texttt{Lookup} в установившемся потоковом профиле
нагрузки. Рассмотрим режим, в котором:

\begin{itemize}
\item в системе хранится таблица из $n$ префиксов
($n \in \{10^4;\,10^5\}$);
\item за единицу времени происходит $\lambda$ инкрементальных
обновлений (Upsert/Delete);
\item клиенты выполняют $\mu$ запросов \texttt{Lookup} в секунду;
\item операция полного \texttt{Dump} вызывается с частотой $\nu$ в
секунду.
\end{itemize}

Тогда суммарную CPU-нагрузку на хранилище можно оценить как

\[
C(\lambda,\mu,\nu) =
   \lambda \cdot t_{\mathrm{upsert}}(n)
 + \mu    \cdot t_{\mathrm{lookup}}(n)
 + \nu    \cdot t_{\mathrm{dump}}(n),
\]

где $t_{\mathrm{upsert}}$, $t_{\mathrm{lookup}}$ и
$t_{\mathrm{dump}}$ -- это средние времена соответствующих операций,
взятые из табл.~\ref{tab:bench-dump}--\ref{tab:bench-shards}.

Подставим типичные значения для $n = 10^5$:

\begin{itemize}
\item \texttt{ranger}: $t_{\mathrm{upsert}} \approx 3{,}94 \cdot 10^{-6}$ с;
$t_{\mathrm{lookup}} \approx 1{,}5 \cdot 10^{-7}$ с;
$t_{\mathrm{dump}} \approx 4{,}89 \cdot 10^{-2}$ с;
\item шард. ART: $t_{\mathrm{upsert}} \approx 1{,}7 \cdot 10^{-7}$ с;
$t_{\mathrm{lookup}} \approx 1{,}5 \cdot 10^{-7}$ с;
$t_{\mathrm{dump}} \approx 1{,}43 \cdot 10^{-1}$ с.
\end{itemize}

Сравнение совокупной стоимости для двух представительных профилей
нагрузки приведено в табл.~\ref{tab:bench-cost}.

\begin{table}[ht]
\centering
\caption{Оценка совокупной CPU-стоимости в установившемся режиме при
$n = 10^5$ (приближённо, CPU-с/с $=$ доля одного CPU-ядра на каждую
секунду физического времени)}
\label{tab:bench-cost}
\small
\begin{tabular}{lrr}
\toprule
Сценарий & \texttt{ranger} & шард. ART \\
\midrule
$\lambda{=}10$, $\mu{=}10^4$, $\nu{=}10^{-2}$
	   & $\approx 0{,}0020$ CPU-с/с
	   & $\approx 0{,}0029$ CPU-с/с \\
$\lambda{=}10^3$, $\mu{=}10^5$, $\nu{=}10^{-2}$
	   & $\approx 0{,}019$ CPU-с/с
	   & $\approx 0{,}017$ CPU-с/с \\
\bottomrule
\end{tabular}
\end{table}

Числа в табл.~\ref{tab:bench-cost} имеют ясную интерпретацию. Оба
варианта остаются на уровне малой доли одного ядра для рассмотренных
профилей, но структура затрат различается: у ART дешевле обновления, а
у \texttt{ranger} дешевле полный \texttt{Dump}. Поэтому профиль нагрузки
определяет выбор backend-а и необходимость кэширования полной выгрузки.

Более того, в этой модели хорошо видно, что доминирующая статья расходов
меняется в зависимости от режима. При частых изменениях важна стоимость
\texttt{Upsert}/\texttt{Delete}; при частых полных выгрузках начинает
доминировать \texttt{Dump} и построение flattened-представления.

Если переложить те же цифры в плоскость требований к задержке
актуализации, важным становится время применения отдельного события.
Пусть допустимая задержка актуализации таблицы
$T_{\mathrm{stale}} \le 1$~с. Потоковая модель достигает этого за счёт
обработки отдельных событий: каждое обновление применяется за единицы
микросекунд или быстрее, а полная выгрузка остаётся отдельным
ресурсно-ёмким сценарием.

%==============================================================
